% Plantilla para articulo en ingles con comandos personalizados David Davalos.
\documentclass[letterpaper,12pt]{article} 

% En Mac modernas, UTF-8 ya es nativo. Evitamos ucs y utf8x por incompatibilidades.
\usepackage[T1]{fontenc} 
\usepackage[spanish]{babel}
\spanishdecimal{.}

\usepackage{color}                    
\usepackage{graphicx} 
\usepackage{url}
\usepackage{bm} % Greek bf
\usepackage{anysize}
\usepackage{amsmath}
\usepackage{amssymb}
\usepackage{multicol}  
\usepackage{texdraw}  
\usepackage{float}
\usepackage{fancyhdr}
\usepackage{amsfonts}
\usepackage{cancel}
\newcommand{\mcE}{\ensuremath{\mathcal{E}}} 
\newcommand{\mcL}{\ensuremath{\mathcal{L}}} 
% NOTA: Si usas pdflatex (por defecto en Mac), pstricks romperá la compilación.
% Si necesitas pstricks, compila usando el motor XeLaTeX.
%\usepackage[usenames,dvipsnames]{pstricks}
%\pst-grad} % For gradients
%\usepackage{pst-plot}

\usepackage{hyperref}
%%%%% De la tesis %%%%%%%%%%
\usepackage[sc]{mathpazo} % ¡Esta es la fuente chida!
\usepackage{lettrine}
\usepackage{amsthm}
\usepackage{xcolor}
\usepackage{tikz}
%%%%%%%%%%%%%%%%%%%%%%%%%%%%%

% Asegúrate de que el archivo 'definiciones.tex' esté en la misma carpeta
%Si no esta usando el iop package
%%Si no esta usando el iop package
%%Si no esta usando el iop package
%\input{definiciones}
%Comandos mios generales
\newcommand{\fig}[1]{fig.\ref{#1}}
\newcommand{\eqn}[1]{eqn.\eqref{#1}}
\newcommand{\tab}[1]{table \ref{#1}}
\newcommand{\medline}{\begin{center}\line(1,0){470}\end{center}}
\newcommand{\suspensivos}{{\color{blue} (.....)}}
\newcommand{\anotacion}[1]{{\color{red} (.. #1 ...)}}
\newcommand{\ie}{\textit{i. e.}}
%Quantum Mechanics Commands
\newcommand{\ket}[1]{{\vert #1 \rangle}}
\newcommand{\bra}[1]{{\langle #1 \vert}}
\newcommand{\proj}[2]{{\vert #1 \rangle \langle #2 \vert}}
\newcommand{\NoM}{{\cal N}}
%Lugares
\newcommand{\unam}{Universidad Nacional Aut\'onoma de M\'exico, M\'exico, D.F., M\'exico}
\newcommand{\icf}{Instituto de Ciencias F\'{\i}sicas, \unam}
\newcommand{\ifunam}{Instituto de F\'{\i}sica, \unam}
\newcommand{\ifunamint}{Instituto de F\'{\i}´ısica, Universidad Nacional Aut\'onoma de M\'exico, M\'exico D.F. 01000, M\'exico}
\newcommand{\fcen}{Departamento de F\'isica ``J. J. Giambiagi'', FCEN, Universidad de Buenos Aires, 1428 Buenos Aires, Argentina}
%%%%%%%%%%% Comandos de Carlos %%%%%%%%%%%%%%%%%%%%%%%%
% enges
%%%%%%%%%%%%%
%%%%%%%%%%%
\newif\ifenglish
\newif\ifespanol
\newif\ifshort
\newif\iflong
\shorttrue
\englishtrue
\newcommand{\enges}[2]{\ifenglish #1 \fi \ifespanol #2 \fi}
%%%%%%
%%%%%

%Comandos mios generales
\newcommand{\fig}[1]{fig.\ref{#1}}
\newcommand{\eqn}[1]{eqn.\eqref{#1}}
\newcommand{\tab}[1]{table \ref{#1}}
\newcommand{\medline}{\begin{center}\line(1,0){470}\end{center}}
\newcommand{\suspensivos}{{\color{blue} (.....)}}
\newcommand{\anotacion}[1]{{\color{red} (.. #1 ...)}}
\newcommand{\ie}{\textit{i. e.}}
%Quantum Mechanics Commands
\newcommand{\ket}[1]{{\vert #1 \rangle}}
\newcommand{\bra}[1]{{\langle #1 \vert}}
\newcommand{\proj}[2]{{\vert #1 \rangle \langle #2 \vert}}
\newcommand{\NoM}{{\cal N}}
%Lugares
\newcommand{\unam}{Universidad Nacional Aut\'onoma de M\'exico, M\'exico, D.F., M\'exico}
\newcommand{\icf}{Instituto de Ciencias F\'{\i}sicas, \unam}
\newcommand{\ifunam}{Instituto de F\'{\i}sica, \unam}
\newcommand{\ifunamint}{Instituto de F\'{\i}´ısica, Universidad Nacional Aut\'onoma de M\'exico, M\'exico D.F. 01000, M\'exico}
\newcommand{\fcen}{Departamento de F\'isica ``J. J. Giambiagi'', FCEN, Universidad de Buenos Aires, 1428 Buenos Aires, Argentina}
%%%%%%%%%%% Comandos de Carlos %%%%%%%%%%%%%%%%%%%%%%%%
% enges
%%%%%%%%%%%%%
%%%%%%%%%%%
\newif\ifenglish
\newif\ifespanol
\newif\ifshort
\newif\iflong
\shorttrue
\englishtrue
\newcommand{\enges}[2]{\ifenglish #1 \fi \ifespanol #2 \fi}
%%%%%%
%%%%%

%Comandos mios generales
\newcommand{\fig}[1]{fig.\ref{#1}}
\newcommand{\eqn}[1]{eqn.\eqref{#1}}
\newcommand{\tab}[1]{table \ref{#1}}
\newcommand{\medline}{\begin{center}\line(1,0){470}\end{center}}
\newcommand{\suspensivos}{{\color{blue} (.....)}}
\newcommand{\anotacion}[1]{{\color{red} (.. #1 ...)}}
\newcommand{\ie}{\textit{i. e.}}
%Quantum Mechanics Commands
\newcommand{\ket}[1]{{\vert #1 \rangle}}
\newcommand{\bra}[1]{{\langle #1 \vert}}
\newcommand{\proj}[2]{{\vert #1 \rangle \langle #2 \vert}}
\newcommand{\NoM}{{\cal N}}
%Lugares
\newcommand{\unam}{Universidad Nacional Aut\'onoma de M\'exico, M\'exico, D.F., M\'exico}
\newcommand{\icf}{Instituto de Ciencias F\'{\i}sicas, \unam}
\newcommand{\ifunam}{Instituto de F\'{\i}sica, \unam}
\newcommand{\ifunamint}{Instituto de F\'{\i}´ısica, Universidad Nacional Aut\'onoma de M\'exico, M\'exico D.F. 01000, M\'exico}
\newcommand{\fcen}{Departamento de F\'isica ``J. J. Giambiagi'', FCEN, Universidad de Buenos Aires, 1428 Buenos Aires, Argentina}
%%%%%%%%%%% Comandos de Carlos %%%%%%%%%%%%%%%%%%%%%%%%
% enges
%%%%%%%%%%%%%
%%%%%%%%%%%
\newif\ifenglish
\newif\ifespanol
\newif\ifshort
\newif\iflong
\shorttrue
\englishtrue
\newcommand{\enges}[2]{\ifenglish #1 \fi \ifespanol #2 \fi}
%%%%%%
%%%%%


% comandos locales
\marginsize{1.5 cm}{1.5 cm}{-0.5 cm}{1 cm}
\author{David Dávalos}                                 
\date{\today}                                 
\title{Estándarización de notación}     

\begin{document} 
\maketitle
Estilo de bibligrafia que funciona bien para todos: \cite{zimansbook}
\begin{itemize}
\item Kets y bras: $\ket{\Psi}$  $\bra{\Psi}$ (en un documento poco denso en vectores de estado, recomiendo usar griegas mayusculas. Si el documento es pesado en vectores de estado, minusculas están bien, como $\ket{\psi}$).
\item Matrices de densidad: $\rho$, $\psi$, $\sigma$.
\item Observables: $A$, $B$, $G$, etc. Con $E$ reservado para POVMs.
\item Espacios de Hilbert: ${\sf H}, {\sf M}$.
\item Espacios de operadores: $\mathcal{B}(\sf H)$ (operadores acotados) o $\mathcal{L}(\sf H)$ (operadores lineales (para dimensión finita son lo mismo que acotados)), $\mathcal{T}(\sf H)$ (operadores trace-class), $\mathcal{S}(\sf H)$ (matrices de densidad).
\item Se pueden usar definiciones de latex, no trabajen doble (en este punto ver el código fuente de este documento.):
% ensuremath hace que jale bien dentro y fuera de ambientes matemáticos
\newcommand{\mcT}{\ensuremath{\mathcal{T}}} %usualmente uno los pone arriba junto a los paquetes.
$\mcT$, \mcT{} (esto fue escrito con un shortcut)
\item Canales cuánticos: $\mathcal{E}$, $\mathcal{M}$ etc, y con corchetes $\mathcal{E}[\psi]$, $\mathcal{E}[\sigma]$, etc. $\mathcal{I}$ suele reservarse para instrumentos cuánticos (que no preservan la traza).
\item Vectores de Bloch y en general vectores de promedios: $\vec r$.
\item Todo promedio/valor esperado debe escribirse como $\langle A \rangle$. Dependiendo del contexto y la necesidad, pueden agregar como subindice el estado respecto al cual se hace el promedio.
\item Matriz de Choi de un canal $\mcE$: $\tau_\mcE$ (observese que es compatible con la notación de matrices de densidad). Entonces, reserven $\tau$ para matrices de Choi (por defecto usen la versión normalizada, que se llama también \textit{estado de Jamio\l kowski}). Entonces $\tau_\mcE=(\text{id}_n \otimes \mcE)[\omega]$, con $\omega=\proj{\Omega}{\Omega}$ (donde ya deben saber que es $\ket{\Omega}$). Usualmente, en el contexto de matrices de Choi $n$ se entiende como la dimensión del espacio de Hilbert del sistema. Adapten eso a sus trabajos.
\item Los regularmente llamados \textit{superoperadores}, como los canales o generadores tipo Lindblad, todos ellos están al mismo ``nivel'', entonces, para generadores de Lindblad se usa la misma notación: $\mcL[\rho]$. Úsense subindices a su criterio, por ejemplo, si $L$ está definido con operadores o parámetros que lo etiqueten $\mcL_{H, \beta}$.
\end{itemize}


% Tu contenido va aquí

\medline % Cambiado de \medline por seguridad, a menos que esté en definiciones.tex
\bibliographystyle{unsrt}
\bibliography{labibliografia}
\end{document}
